\documentclass[10pt,a4paper]{article}
\usepackage[utf8]{inputenc}
\usepackage{amsmath}
\usepackage{mathtools}
\usepackage{amsfonts}
\usepackage{amssymb}
\usepackage{enumerate}
\usepackage{hyperref}
\usepackage{graphicx}
%\usepackage{natbib}
\bibliographystyle{apalike}
%\setcitestyle{authoryear,open={((},close={))}}

\newcommand{\delims}[3]{\!\left #1 #2 \right #3}
\newcommand{\parens}[1]{\delims{(}{#1}{)}}
\newcommand{\braces}[1]{\delims{\lbrace}{#1}{\rbrace}}
\newcommand{\brackets}[1]{\delims{[}{#1}{]}}
\newcommand{\norm}[1]{\delims{\|}{#1}{\|}}
\newcommand{\abs}[1]{\delims{|}{#1}{|}}
\newcommand{\inner}[2]{\delims{\langle}{#1,#2}{\rangle}}
\newcommand{\vectinner}[2]{\inner{\vect{#1}}{\vect{#2}}}

\newcommand{\vect}[1]{\mathbf{#1}}

\begin{document}
\title{\vspace{-10ex}Top-down and Bottom-up Processing}
\author{Dmitry Brizhinev}
\maketitle

\section{Introduction}
I have found that the concept of top-down vs bottom-up processing is widely used in psychology and neuroscience \cite{citationneeded}. Most visual information from human eyes passes via the lateral geniculate nucleus. This is where the first post-retinal processing happens to visual information \cite{citationneeded}. However, most of the \emph{inputs} to the LGN are from higher brain regions \cite{citationneeded}. This suggests that even early processing depends a great deal on top-down cues.

\section{Definition}
I will explain what I mean by top-down and bottom-up processing using an example. Consider this image \cite{citationneeded}:

\begin{figure}[h!]
  \includegraphics[width=200pt]{thecat.png}
\end{figure}

A human can effortlessly identify this as the text \texttt{THE CAT}. This is despite the fact that the \texttt{H} and \texttt{A} glyphs are identical. Any system that tries to read the text by \emph{first} identifying the letters via combinations of pixels (or strokes), and \emph{then} identifying the words from the letters, is doomed to fail. Such a system can only produce the outputs \texttt{THE CHT} or \texttt{TAE CAT}. [This would be a good place for an empirical demonstration, or for citing one \cite{researchneeded}] To correctly read this text, a system must use the context of the glyph to help identify the glyph. We refer to something like \emph{looking at the pixels in an image to identify the letter} as bottom-up processing, and \emph{looking at the letters around it to figure out what the letter \textbf{should be}} as top-down processing.

A naive model of cognition involves only bottom-up processing. Such models typically define a series of processing steps, each converting a simpler, less meaningful representation into a more complex one, e.g. \emph{pixels $\rightarrow$ strokes $\rightarrow$ letters $\rightarrow$ words $\rightarrow$ sentences}. Top-down processing is any part of the model where information from a ``later" stage in this chain is used to help with an earlier stage. This is possible after the system has made a first guess at the later stage using only bottom-up information, but it might also involve \emph{what the system expected to see}, before any bottom-up information is available.

This entire discussion hinges on the hypothesis that it is possible to clearly delineate ``higher" brain regions, and ``later" stages of processing, \emph{even though} information can flow in any direction. For the early visual system, this is possible -- we can just count the minimum number of synapses it takes for information to travel to a region from the retina. However, perhaps in the frontal cortex these distinctions cease to make sense.

Side note: I am not sure what the relationship is between \emph{top-down/bottom-up processing} as I have defined it, and the \emph{global/local processing} dichotomy also talked about in psychology \cite{citationneeded}.

\section{Approaches to top-down vs bottom-up}
\subsection{Neuroscience and psychology}
The concept of top-down and bottom-up processing originates from psychology. In psychology, top-down is defined something like my definition above \cite{citationneeded}. In neuroscience, top-down and bottom-up refer to the direction of information flow (neural connections) \cite{citationneeded}. It is clear that human cognition involves both kinds of processing, and that top-down processing is very important (rather than a minor part) \cite{citationneeded}.

\subsection{Coward}
The feature extraction component of Coward's architecture implements a particular kind of top-down processing \cite[p.~??]{coward}\cite{citationneeded}. Coward imagines that the feature extraction modules receive information from simpler modules and process it to detect (or not) their feature. This is a form of bottom-up processing. Top-down processing is implemented via what Coward calls \emph{indirect} activation -- one of the actions that the architecture can take in response to inputs is to alter the activation of some feature extractors (a simple example would be the \emph{action} of recalling a memory). Coward posits that the feature extraction modules would have two dedicated sets of inputs -- \emph{direct} (bottom-up) inputs from other feature extractors, and \emph{indirect} (top-down) inputs from the action recommendation component that forcibly activate the feature extractor if the action recommendation component deems it necessary. Coward imagines that there would be inhibitory connections between these two inputs, so that a feature extractor can experience direct \emph{or} indirect activation, but not both. Coward believes this is necessary to avoid hallucinations, where an indirect activation (like memory retrieval) is confused with a direct one (like vision).

I would need to recheck to see whether there is any top-down/bottom-up distinction in the action recommendation part of his model \cite{researchneeded}.

\subsection{Kurzweil}
Kurzweil's architecture proposal \cite{kurzweil} is much less detailed than Coward's, but it pays special attention to top-down processing, using as an example the case of identifying a partially-obscured or corrupted letter in a word \cite[p.~??]{kurzweil}\cite{citationneeded}. Unlike Coward, Kurzweil believes human brain architecture must combine bottom-up and top-down cues, to allow the identification of a letter where neither bit alone is enough. His proposed modules are very similar to Coward's feature extractors, but they do not have a dedicated inhibition circuit to prevent mixing top-down and bottom-up information.

Consider the word \texttt{M??IMUM}. It could be either \texttt{MAXIMUM} or \texttt{MINIMUM}, so top-down processing is not enough to fill in the missing letters (unless the whole word can be inferred from context). Suppose also the missing letters are partially obscured, so bottom-up processing can only produce an ambiguity between \texttt{AX} and \texttt{AY}. Combining both top-down and bottom-up information removes the ambiguities.

\subsection{Minsky}
From what I remember, Minsky's model is a (tree-like) hierarchy of behavioural modules. Thus, I don't think it handles top-down processing very well. However, I will have to go back over the text to make sure \cite{citationneeded}.

\subsection{CNN}
The standard convolutional neural network is strictly bottom-up. In general, any neural net that is based on only one time step of processing will be limited, because it cannot use the results of one time step to make a better decision the second time around (e.g. use the first guess to bias the lower-level neurons in particular ways).

The approach of adding \emph{skip} links joining early layers directly to high level ones \cite{citationneeded} might crudely reproduce top-down processing. If a late layer has access to both high-level features and the raw pixels, it might use the features to re-do the low-level processing. If a network actually learns this function, it will presumably be more efficient to reduce the number of layers and replace it with a different architecture that supports top-down processing explicitly.

Using backpropagation to train a neural network is a form of top-down information flow in a trivial sense -- information flows from \emph{later} layers to \emph{earlier} ones. However, I feel like this doesn't capture the spirit of the idea -- that we can improve our performance at test time (rather than training time) by using the outputs of a later layer to tune an earlier one.

\subsection{RNN}
The recurrent neural network potentially supports top-down processing of some kind, since top-level outputs are fed back in at test time \cite{citationneeded}. If we give the network an input, then wait several steps before expecting an output, the network should in-principle be able to learn the strategy of using processed results to bias lower-level features and produce a better result.

I should look more into some specific RNNs, particularly applications of RNNs to combined optical-character-recognition plus higher processing tasks (e.g. text-to-speech). \cite{researchneeded}

Side note: I find it interesting that neuroscience tends to use the idea of \emph{neural connections going backwards and forwards}, while the AI field has settled on the metaphor of \emph{network outputs feeding into its inputs}, and of \emph{unrolling the RNN into a sequential network with periodic constraints on the weights}.

\subsection{AIXI}
AIXI does not have the concept of \emph{brain regions} or \emph{processing steps}, and so the distinction between top-down and bottom-up processing does not apply to it \cite{citationneeded}. I think of this as an example of the general reason why AIXI is insufficient as a model of \emph{cognition} (not that it's intended to be) -- it defines only the output of the process, not how a realistic system could arrive at the output.

I also want to look at approximate AIXI implementations, to see whether the situation is different there. \cite{researchneeded}



\subsection{Watson}
I don't know enough about the architecture to answer this yet \cite{researchneeded}.

\subsection{AlphaGo}
I don't know enough about the architecture to answer this yet \cite{researchneeded}.

\subsection{Deep Q-learning}
I don't know enough about the architecture to answer this yet \cite{researchneeded}.

%\cite[p.~150--153]{ref} \cite[p.~150]{ref}.





\clearpage

\bibliography{../references}
\end{document}